% Fibonacci-Pell nearest gaps: an all-exponent classification and
% quadratic-unit orbit rigidity
%
% Authors:
%   Dr. Denys Dutykh (Mathematics Department, Khalifa University of Science
%   and Technology, Abu Dhabi, UAE)
%   Prof. Laurent Vuillon (Univ. Savoie Mont Blanc, CNRS, LAMA, Chambery,
%   France)
%
\section{The odd-exponent branch}
\label{sec:odd-exponents}

The preceding sections classified the even-exponent branch. We now return to
the intrinsic setup of \cref{sec:system} and suppose that the exponent $n$ is
odd. The parity selector forces $3\mid q$, so this branch is disjoint from
the primitive-rectangle argument. Its content has a different factorisation,
and in particular no analogue of the divisor-height separation
\cref{lem:height-separation} is available. The proof closes after reversing
the order of the two logarithmic-form estimates.

\begin{proposition}[Negative unit exponent in the odd branch]
\label{prop:odd-negative-exponent}
Suppose that $q\geqslant1$, $\sigma\in\{1,-1\}$, and a positive odd integer
$n$ satisfy \cref{eq:nearest-window,eq:norm-content}. Then $3\mid q$, and
the unit normalisation has the form

\begin{equation}
 \lambda^n-D_{q,\sigma}^2=g\lambda^{-h},
 \qquad
 h\geqslant3\ \text{odd}.
 \label{eq:odd-negative-target}
\end{equation}
\end{proposition}

\begin{proof}
By \cref{prop:parity-selector}, $3\mid q$ and hence $q\geqslant3$. Put

\begin{equation}
 A_+=F_q+F_{q+1}\sqrt2,
 \qquad
 A_-=-F_q+F_{q+1}\sqrt2.
 \label{eq:odd-Apm}
\end{equation}

Both are positive. By \cref{lem:normalisation}, there is a unique odd
integer $j$ such that

\begin{equation}
 \lambda^n-D_{q,\sigma}^2=g\lambda^j.
 \label{eq:odd-unit-j}
\end{equation}

Since $\overline{D_{q,\sigma}}=-D_{q,-\sigma}$ and both $n$ and $j$ are
odd, conjugation gives

\begin{equation}
 \lambda^{-n}+D_{q,-\sigma}^2=g\lambda^{-j}.
 \label{eq:odd-unit-conjugate}
\end{equation}

Division of \cref{eq:odd-unit-j} by
\cref{eq:odd-unit-conjugate}, followed by the strict upper window, yields

\begin{equation}
 \lambda^{2j}
 <\frac{D_{q,\sigma}^2}{D_{q,-\sigma}^2}.
 \label{eq:odd-j-ratio}
\end{equation}

The ratio bounds \cref{lem:ratio-bounds} say
$\lambda A_-<A_+<3A_-$. Thus $j\leqslant1$ for $\sigma=1$, whereas
$j\leqslant-3$ for $\sigma=-1$.

It remains to remove $j=1$ and $j=-1$ on the plus branch. The following
coefficient argument in fact excludes those values for either sign. Set

\begin{equation}
 H_+=F_{q+1}^2+F_{q+2}^2=F_{2q+3},
 \qquad
 H_-=F_{q+1}^2+F_{q-1}^2
 =3F_q^2+2(-1)^q.
 \label{eq:odd-seed-heights}
\end{equation}

If $j=1$, multiply \cref{eq:odd-unit-j} by $\lambda^{-1}$ and compare
the $\sqrt2$-coefficients; if $j=-1$, multiply it by $\lambda$ instead.
The resulting equations are

\begin{equation}
 P_{n-1}=H_{-\sigma}
 \quad\text{or}\quad
 P_{n+1}=H_\sigma,
 \label{eq:odd-leading-intersections}
\end{equation}

respectively. Alekseyev's complete Fibonacci--Pell intersection theorem
states that the common values of the two sequences are $0,1,2,5$
\cite[Theorem~9]{Alekseyev2011}. Hence $H_+=F_{2q+3}\geqslant F_9=34$
cannot occur. Alekseyev and Tengely's complete near-square table gives

\begin{equation}
 3y^2+2\in\{2,5,29\},
 \qquad
 3y^2-2\in\{1\}
 \label{eq:odd-near-square-values}
\end{equation}

whenever the displayed number is a Pell value
\cite[Table~1]{AlekseyevTengely2014}. If $q$ is odd, then
$H_-=3F_q^2-2\geqslant10$; if $q$ is even, then
$H_-=3F_q^2+2\geqslant194$. Thus neither alternative in
\cref{eq:odd-leading-intersections} is possible. Every remaining $j$ is at
most $-3$, proving \cref{eq:odd-negative-target} with $j=-h$.
\end{proof}

\begin{lemma}[Odd trace and content factorisation]
\label{lem:odd-factorisation}
Under the hypotheses of \cref{prop:odd-negative-exponent}, put

\begin{equation}
 f=F_q,
 \qquad
 x=F_{q+1},
 \qquad
 L=f^2+2x^2,
 \qquad
 R_q=2x^2-f^2,
 \qquad
 M=n+h.
\end{equation}

Then

\begin{equation}
 \lambda^n-\lambda^{-n-2h}
 =D_{q,\sigma}^2+D_{q,-\sigma}^2\lambda^{-2h},
 \label{eq:odd-trace}
\end{equation}

and

\begin{equation}
 R_q^2+1=g(2C_M-g),
 \qquad
 1\leqslant g<R_q,
 \qquad
 g\mid R_q^2+1.
 \label{eq:odd-content-factorisation}
\end{equation}

In particular, every prime divisor of $g$ is congruent to $1$ modulo $4$.
\end{lemma}

\begin{proof}
Conjugating \cref{eq:odd-negative-target} gives

\begin{equation}
 \lambda^{-n}+D_{q,-\sigma}^2=g\lambda^h.
 \label{eq:odd-negative-conjugate}
\end{equation}

Eliminating $g$ between this equation and
\cref{eq:odd-negative-target} proves \cref{eq:odd-trace}. Since
$\lambda^{-h}=-C_h+P_h\sqrt2$, coefficient comparison gives

\begin{equation}
 L-C_n=gC_h,
 \qquad
 P_n-2\sigma fx=gP_h.
 \label{eq:odd-coefficient-comparison}
\end{equation}

Taking norms in \cref{eq:odd-negative-target}, and using that $M=n+h$ is
even, gives

\begin{equation}
 R_q^2+1=g(2C_M-g).
\end{equation}

The Pell addition formulas and \cref{eq:odd-coefficient-comparison} also
give

\begin{equation}
 C_M-g=LC_h+4\sigma fxP_h.
 \label{eq:odd-positive-factor}
\end{equation}

This is positive for $\sigma=1$. For $\sigma=-1$, use

\begin{equation}
 L-3fx=(x-f)(2x-f)>0,
 \qquad
 3C_h>4P_h
\end{equation}

to obtain positivity again. Hence $g<C_M$, so
$g^2<R_q^2+1$. Here $g$ is odd by the coefficient parity in
\cref{eq:gap-coefficients}, whereas $R_q$ is even because $3\mid q$.
Therefore $g<R_q$, and the divisibility assertion follows from the first
factorisation in \cref{eq:odd-content-factorisation}. Finally, if a prime
$p$ divides $g$, then $p$ is odd and $R_q^2\equiv-1\pmod p$; hence
$p\equiv1\pmod4$.
\end{proof}

The factorisation in \cref{lem:odd-factorisation} explains why the even
proof cannot simply be quoted. It gives no lower bound for $g$ and no
inequality $h<n$. We therefore use the fixed logarithmic form first, control
$\min\{h,q\gamma\}$, and only then invoke the moving form.

\begin{theorem}[Odd-exponent exclusion]
\label{thm:odd-exclusion}
For $q\geqslant1$ and $\sigma\in\{1,-1\}$, no positive odd integer $n$
can satisfy both \cref{eq:nearest-window,eq:norm-content}.
\end{theorem}

\begin{proof}
Suppose otherwise. By \cref{prop:odd-negative-exponent}, $3\mid q$ and
\cref{eq:odd-negative-target} holds. We first treat $q\geqslant7$; the
indices $q=3,6$ will be included in the exact terminal enumeration.

Retain $\varphi$, $\epsilon=(-1)^q$, and the four Binet coefficients from
\cref{eq:four-coefficients,eq:two-binet}, and put

\begin{equation}
 \gamma=\frac{\log\varphi}{\log\lambda},
 \qquad
 H=1+\log(2q).
 \label{eq:odd-gamma-H}
\end{equation}

The window and the elementary bound $D_{q,\sigma}<4\varphi^q$ give

\begin{equation}
 \lambda^n<32\varphi^{2q}<\lambda^{2q},
 \qquad
 n<2q.
 \label{eq:odd-n-less-2q}
\end{equation}

Indeed, $\lambda/\varphi>7/5$ and $(7/5)^{14}>32$.

We begin with the fixed form

\begin{equation}
 \Lambda_0
 =n\log\lambda-2q\log\varphi
 -\log\left(\frac{u_\sigma^2}{5}\right).
 \label{eq:odd-Lambda0}
\end{equation}

Define

\begin{equation}
 \delta_\sigma
 =\epsilon\frac{u'_\sigma}{u_\sigma}\varphi^{-2q},
 \qquad
 \widetilde\delta_\sigma
 =\epsilon\frac{v'_\sigma}{v_\sigma}\varphi^{-2q}.
\end{equation}

Expanding \cref{eq:odd-trace} gives the exact identity

\begin{align}
 e^{\Lambda_0}-1={}&
 2\delta_\sigma+\delta_\sigma^2
 +\frac{v_\sigma^2}{u_\sigma^2}
 (1+\widetilde\delta_\sigma)^2\lambda^{-2h}\notag\\
 &+\frac5{u_\sigma^2}
 \lambda^{-n-2h}\varphi^{-2q}.
 \label{eq:odd-fixed-expansion}
\end{align}

The final sign is positive; this is the precise change from
\cref{eq:fixed-exact-expansion}. The coefficient estimates used there are
insensitive to this sign and give

\begin{equation}
 0<\abs{\Lambda_0}
 <12\left(\lambda^{-2h}+\varphi^{-2q}\right).
 \label{eq:odd-Lambda0-small}
\end{equation}

Nonvanishing is exact: if $\Lambda_0=0$, full norms in
$\Q(\sqrt2,\sqrt5)$ would give $1=5^{-4}$. By
\cref{lem:matveev,eq:matveev-prefactor,eq:odd-n-less-2q}, the same explicit
height parameters as in \cref{sec:height} yield

\begin{equation}
 \log\abs{\Lambda_0}>
 -1.855\mathbin{\cdot}10^{14}H.
 \label{eq:odd-Lambda0-lower}
\end{equation}

Put

\begin{equation}
 w=\min\{h,q\gamma\}.
\end{equation}

The right-hand side of \cref{eq:odd-Lambda0-small} is at most
$24\lambda^{-2w}$. Since $\log\lambda>4/5$ and $\log24<4$, comparison with
\cref{eq:odd-Lambda0-lower} gives

\begin{equation}
 w<1.16\mathbin{\cdot}10^{14}H.
 \label{eq:odd-w-logarithmic}
\end{equation}

We next use the moving coefficient $\mathcal A_{\sigma,h}$ from
\cref{eq:moving-coefficient}. The form

\begin{equation}
 \Lambda_1
 =n\log\lambda-2q\log\varphi
 -\log\mathcal A_{\sigma,h}
 \label{eq:odd-Lambda1}
\end{equation}

has the exact remainder

\begin{align}
 \mathcal R^{\mathrm{odd}}_{\sigma,q,h}={}&
 \frac{2\epsilon}{5}
 \left((1-\sigma\sqrt2)
 +(1+\sigma\sqrt2)\lambda^{-2h}\right)\notag\\
 &+\frac{\varphi^{-2q}}5
 \left((u'_\sigma)^2
 +(v'_\sigma)^2\lambda^{-2h}\right)
 +\lambda^{-n-2h}.
 \label{eq:odd-moving-remainder}
\end{align}

Thus

\begin{equation}
 \lambda^n
 =\mathcal A_{\sigma,h}\varphi^{2q}
 +\mathcal R^{\mathrm{odd}}_{\sigma,q,h}.
\end{equation}

Only the sign of the last term differs from
\cref{eq:moving-remainder}. The same absolute estimates, full norm
\cref{eq:A-norm}, and height bound \cref{eq:A-height} therefore give

\begin{equation}
 0<\abs{\Lambda_1}<21\varphi^{-2q}
 \label{eq:odd-Lambda1-small}
\end{equation}

and, by \cref{lem:matveev,eq:matveev-prefactor},

\begin{equation}
 2q\log\varphi
 <9.275\mathbin{\cdot}10^{12}(4h+44)H+\log21.
 \label{eq:odd-q-h-comparison}
\end{equation}

If $q\geqslant10^{32}$, then $\gamma>1/2$ and the monotonicity of $q/H$
show that $q\gamma$ exceeds the right-hand side of
\cref{eq:odd-w-logarithmic}. Hence $w=h$. At $q_0=10^{32}$ one has
$H(q_0)<79$. Substitution of
$h<1.16\mathbin{\cdot}10^{14}H$ into
\cref{eq:odd-q-h-comparison}, using $\log\varphi>2/5$, gives at $q_0$

\begin{equation}
 q<\frac54\left[
 9.275\mathbin{\cdot}10^{12}
 \left(4(1.16\mathbin{\cdot}10^{14})H+44\right)H+4
 \right]
 <3.36\mathbin{\cdot}10^{31}.
 \label{eq:odd-crude-contradiction}
\end{equation}

After division by $q$, the terms $H^2/q$, $H/q$, and $1/q$ decrease for
$q\geqslant q_0$, so the contradiction persists. Consequently

\begin{equation}
 q<10^{32}.
 \label{eq:odd-q-crude}
\end{equation}

It remains to apply the exact reductions already used in
\cref{sec:classification}. Put

\begin{equation}
 \mu_\sigma
 =\frac{\log(u_\sigma^2/5)}{\log\lambda}.
\end{equation}

Dividing \cref{eq:odd-Lambda0-small} by $\log\lambda$ gives

\begin{equation}
 0<\abs{2q\gamma-n+\mu_\sigma}
 <30(\lambda^2)^{-w},
 \qquad
 2q<2\mathbin{\cdot}10^{32}.
 \label{eq:odd-fixed-reduction}
\end{equation}

The common rational approximation and the two fixed-shift certificates
\cref{eq:common-convergent,eq:convergent-certificate,eq:fixed-eta,eq:fixed-cutoff}
depend only on $\gamma$, $\mu_\sigma$, and $\lambda$, not on the parity of
$n$. The absolute-value reduction in \cref{lem:absolute-reduction}
therefore proves

\begin{equation}
 w<48.
 \label{eq:odd-w-48}
\end{equation}

If $q\geqslant96$, then $q\gamma>48$, so $h=w<48$. Put

\begin{equation}
 \nu_{\sigma,h}
 =\frac{\log\mathcal A_{\sigma,h}}{\log\lambda}.
\end{equation}

\Cref{eq:odd-Lambda1-small} gives

\begin{equation}
 0<\abs{2q\gamma-n+\nu_{\sigma,h}}
 <27(\varphi^2)^{-q}.
 \label{eq:odd-moving-reduction}
\end{equation}

The moving certificates \cref{eq:moving-eta,eq:moving-cutoff} cover both
signs and every odd $3\leqslant h<192$, hence in particular the present
range $3\leqslant h<48$. A second application of
\cref{lem:absolute-reduction} gives $q<90$, contradicting
$q\geqslant96$. We conclude that

\begin{equation}
 q<96.
 \label{eq:odd-q-96}
\end{equation}

The remaining indices are exactly

\begin{equation}
 q=3,6,\ldots,93,
 \qquad
 \sigma\in\{1,-1\}.
 \label{eq:odd-terminal-domain}
\end{equation}

Because $\lambda>2$, each strict factor-two window contains at most one
integral power. Exact recurrence arithmetic constructs every candidate in
\cref{eq:odd-terminal-domain}, computes its coefficient content, and tests
the norm identity. The complete terminal list is empty. This contradiction
proves the theorem.
\end{proof}

Together with the even branch and the direct $q=1$ check in
\cref{eq:q1-boundary}, the odd exclusion proves
\cref{thm:classification-intro}.

The order of the argument is essential. The fixed Matveev estimate first
controls $\min\{h,q\gamma\}$; only in the large-$q$ alternative does this
become a bound for $h$. The moving estimate then bounds $q$, after which the
same fixed and moving rational-interval certificates reduce the problem to
the finite domain \cref{eq:odd-terminal-domain}. No finite search is
extrapolated beyond its proved range.

The all-exponent nearest-gap classification is now complete. We next
reinterpret it dynamically, where it becomes the arithmetic input to a
complete Pell-orbit theorem.
